\documentclass[book.tex]{subfiles}

\begin{document}
\chapter{Lie Groups and Lie Algebra}
\label{chap:lie}
\noindent\rule{11cm}{0.4pt} \\
\textbf{Prerequisites:} \autoref{chap:math_basics}, \autoref{chap:manifold} \\ 
\textbf{Difficulty Level:} *** \\
\noindent\rule{11cm}{0.4pt}
\vspace{5mm}

Lie groups (pronounced "Lee") are mathematical groups that are also differentiable manifolds. Lie groups arise repeatedly in quantum mechanics as the correct structure to model continuous groups of transformations. As we will see in this chapter, the group of rotations in three dimensions is given by $\mathrm{SO}(3)$, a Lie group. Or the group of distance preserving transformations is given by $\mathrm{E}(3)$. Another foundational concept is that of the representation of Lie groups

\section{Definitions}

A real Lie group is a group which is also a finite-dimensional real smooth manifold. The group is written multiplicatively, and the group operations of multiplication and inversion are smooth maps on the manifold. Formally, the map 
\begin{align}
    \mu:G \times G \to G \quad \mu(x, y) = xy
\end{align}
is a smooth map from the product manifold $G \times G$ to $G$. %(\textcolor{red}{TODO: Define a product manifold in manifold chapter.})

Consider the set of $2 \times 2$ invertible matrices
\begin{align}
    \textrm{GL}(2, \mathbb{R}) &= \left \{A = \begin{pmatrix} a & b \\ c & d \end{pmatrix} | \det A \neq 0\right \} 
\end{align}
We claim that $\textrm{GL}(2, \mathbb{R})$ is a Lie group. More generally, $\textrm{GL}(n, \mathbb{R})$ of $n \times n$ invertible matrices is a Lie group as is $\textrm{GL}(n, \mathbb{C})$. So too are the special linear groups of matrices with determinant $1$
\begin{align}
    \textrm{SL}(n, \mathbb{R}) &= \{ A \in M_{n\times n}(\mathbb{R}) | \det A = 1 \} \\
    \textrm{SL}(n, \mathbb{C}) &= \{ A \in M_{n\times n}(\mathbb{C}) | \det A = 1 \} \\
\end{align}
The orthogonal and special orthogonal groups of orthogonal real matrices form a Lie group
\begin{align}
    O(n) &= \{ A \in M_{n \times n}(\mathbb{R}) | \det A \neq 0, A^T = A^{-1} \} \\
    SO(n) &= \{ A \in M_{n \times n}(\mathbb{R}) | \det A = 1, A^T = A^{-1} \}
\end{align}
as do the unitary and special unitary groups of complex matrices.
\begin{align}
    U(n) &= \{ U \in M_{n \times n}(\mathbb{C}) | \det U \neq 0, U^* = U^{-1} \} \\
    SU(n) &= \{ U \in M_{n \times n}(\mathbb{C}) | \det U = 1, U^T = U^{-1} \}
\end{align}

\section{Lie Algebras}

Every Lie group has associated with it the associated Lie algebra. 

%\textcolor{red}{TODO: Define an algebra in the opening chapter.}

\begin{align}
    \textrm{Lie}(G) &= \{ X \in M(n, \mathbb{C}) | \textrm{exp}(t X) \in G, \ \forall t \in \mathbb{R} \}
\end{align}
This algebra has with it an associated bracket structure
\begin{align}
    [X, Y] &= XY - YX
\end{align}

\section{Exercises}
\begin{enumerate}
    \item Prove that $\textrm{GL}(2, \mathbb{R})$ is a Lie group
    \item Prove that $\textrm{GL}(n, \mathbb{R})$ is a Lie group
    \item Prove that $\textrm{Lie}(G)$ is an algebra
\end{enumerate}

\end{document}